%% annex/p06/cardinals-choice.tex --- proof annex for P6 cardinals and choice
%% One annex item per proved=/sketched= result, in body order.

\begin{annexitem}{annex:p6.card.1}{thm:p06-hartogs}{Hartogs's theorem}
\sketchnote{the coding of well-orders on subsets of \(A\), the Replacement
instance collecting their collapse ordinals, and the Mostowski-collapse
uniqueness theorem are imported from \cite{Jech2003}.}
The well-orders whose fields are subsets of \(A\) form a set.  Collapse each
to its unique ordinal and collect the resulting ordinals into a set \(W\).
If \(\gamma=\sup W+1\) injected into \(A\), the injection would transport the
well-order on \(\gamma\) to a well-order on a subset of \(A\), putting an
ordinal of order type \(\gamma\) into \(W\), contrary to
\(\gamma>\sup W\).  Hence some ordinal does not inject into \(A\); the least
such ordinal exists by ordinal well-ordering and is \(h(A)\).
\end{annexitem}

\begin{annexitem}{annex:p6.choice.1}{thm:p06-zorn}{Choice equivalents and Zorn's lemma}
\sketchnote{set-coding and Replacement in the Hartogs-stopping argument;
the reductions for Hausdorff's maximal principle, Tukey's lemma, and
Tychonoff's theorem, imported from \cite[Chapter~5]{Jech2003}; and the cited
calibration of compact Hausdorff Tychonoff against the Boolean Prime Ideal
Theorem.}
Assume Choice.  In a poset \(P\) in which every chain has an upper bound, fix a
well-order of \(P\).  Recursively for \(\alpha<h(P)\), choose the least strict
upper bound of the previously chosen chain whenever one exists.  If the
construction ran through every \(\alpha<h(P)\), its distinct values would
inject \(h(P)\) into \(P\), contradicting Hartogs's theorem.  At the first
stopping stage, the chain hypothesis supplies an upper bound \(u\).  Since no
strict upper bound exists, \(u\) is the greatest member of the chain; and if
some \(u'>u\) existed, \(u'\) would be a strict upper bound.  Thus \(u\) is
maximal, proving Zorn's lemma from Choice.

Apply Zorn's lemma to partial well-orders of an arbitrary set, ordered by
end-extension.  Unions of chains are upper bounds; a maximal partial
well-order must cover the whole set, since otherwise one more point could be
adjoined.  Thus Zorn implies the Well-ordering Theorem.  A well-order of the
union of a set-indexed family of nonempty sets selects the least member of each
factor, producing a point of its product.  Conversely, a point of every such
product is exactly a choice function.  This closes the migrated implication
cycle
\[
  \begin{aligned}
    \mathsf{AC}&\ \text{ implies Zorn, which implies well-ordering;}\\
    \text{well-ordering}&\ \text{ implies the product principle, which implies }
      \mathsf{AC}.
  \end{aligned}
\]
The remaining equivalent rows of Table~\ref{tab:p06-ac-equivalents} are the
standard cited refinements of this cycle.
\end{annexitem}

\begin{annexitem}{annex:p6.choice.2}{thm:p06-konig-infinity}{K\"onig's infinity lemma for the binary tree}
\begin{proof}
Call a finite binary sequence \(\mathbf s\) productive iff
\(\mathbf s\in\mathcal T\) and, for every
\(m\in\N\), it has an extension in \(\mathcal T\) of length
\(\operatorname{len}(\mathbf s)+m\).  The empty sequence is productive because
\(\mathcal T\) has a node at every level.  If a productive \(\mathbf s\) had neither
\(\mathbf s^\frown0\) nor \(\mathbf s^\frown1\) productive, each child would have a finite
bound on the heights of its extensions: this is immediate if the child is not
in \(\mathcal T\), and otherwise follows from the failure of productivity.  The maximum
of those two bounds would bound the height above \(\mathbf s\), a contradiction.
Recursively extend a productive node by \(0\) when its
\(0\)-child is productive, and by \(1\) otherwise.  This is a definition by
the fixed two-element order, not a sequence of arbitrary choices.  The union
of the recursively defined nodes is an \(f\in2^\omega\), and every
\(f\mathord{\upharpoonright}n\) lies in \(\mathcal T\); hence \(f\) is a branch.
\end{proof}
\end{annexitem}

\begin{annexitem}{annex:p6.choice.3}{thm:p06-classical-fan}{Classical fan theorem}
\begin{proof}
Suppose no uniform \(N\) exists.  Let \(\mathcal T\) consist of the binary sequences
none of whose initial segments belongs to \(B\).  This is a binary tree, and
the failure of a uniform bound supplies a node at every level.  By
Theorem~\ref{thm:p06-konig-infinity}, \(\mathcal T\) has a branch \(f\).  No initial
segment of \(f\) belongs to \(B\), contradicting that \(B\) is a bar.
\end{proof}
\end{annexitem}
